\input{configTD}

\usepackage{tikz}

\title{Analyse Numérique TD1 --- Interpolation et intégration numérique}

\author{}
\date{}

\begin{document}
\sffamily
\maketitle

% ====================================================================
%  RAPPELS DE COURS : INTERPOLATION
% ====================================================================

\begin{mdframed}[backgroundcolor=ThemeLight, linecolor=Theme,
  linewidth=1.5pt, innertopmargin=10pt, innerbottommargin=10pt]
\textbf{\large Interpolation polynomiale de Lagrange}

\medskip

\textbf{Problème.}
On dispose de $n+1$ points de mesure $(x_0, y_0),\,(x_1, y_1),\,\ldots,\,(x_n, y_n)$
avec des abscisses~$x_i$ \emph{distinctes}.

On cherche un polynôme $P$ qui passe par tous les points, c'est à dire tel que
\[
  P(x_i) = y_i, \quad i = 0, 1, \ldots, n.
\]

\begin{center}
\begin{tikzpicture}[scale=0.9]
  % Axes
  \draw[->,gray] (-0.3,0) -- (7.2,0) node[right] {\small$x$};
  \draw[->,gray] (0,-0.3) -- (0,3.8) node[above] {\small$y$};

  % Data points
  \coordinate (P0) at (0.5, 1.2);
  \coordinate (P1) at (2.0, 3.0);
  \coordinate (P2) at (3.5, 2.0);
  \coordinate (P3) at (5.0, 3.2);
  \coordinate (P4) at (6.5, 1.5);

  % Interpolation polynomial (smooth curve through points)
  \draw[thick, Theme, smooth, tension=0.5]
    plot coordinates {(0.5,1.2) (2.0,3.0) (3.5,2.0) (5.0,3.2) (6.5,1.5)};

  % Dashed lines to axes
  \foreach \p/\lab/\ylab in {P0/$x_0$/$y_0$, P1/$x_1$/$y_1$,
    P2/$x_2$/$y_2$, P3/$x_3$/$y_3$, P4/$x_4$/$y_4$} {
    \draw[dotted, gray] (\p) -- (\p |- 0,0);
    \node[below, font=\scriptsize] at (\p |- 0,0) {\lab};
  }

  % Points
  \foreach \p in {P0,P1,P2,P3,P4}
    \filldraw[Theme!80!black] (\p) circle (2pt);

  % Label for the curve
  \node[Theme!80!black, font=\small, anchor=south west] at (5.5,2.8) {};

  % Question mark between points
  \draw[<->, Accent, thick] (2.7,2.5) -- (2.7,2.5)
    node[above, font=\small\itshape, Accent!80!black] {$P(x_?)$};
  \filldraw[Accent] (2.7, 2.57) circle (1.8pt);
  \draw[dotted, Accent!70!black] (2.7,0) -- (2.7,2.57);
  \node[below, font=\scriptsize, Accent!80!black] at (2.7,0) {$x_?$};
\end{tikzpicture}
\end{center}

\textbf{Polynômes de base de Lagrange.}
Pour chaque $i \in \{0, \ldots, n\}$, on définit
\[
  \ell_i(x)
  = \prod_{\substack{0 \leq j \leq n \\ j \neq i}}
    \frac{x - x_j}{x_i - x_j}\,.
\]
Ce polynôme de degré $n$ vaut $1$ en $x_i$ et $0$ en chaque autre $x_j$.

\medskip

\textbf{Polynôme d'interpolation.}
Le polynôme
\[
  P(x) = \sum_{i=0}^{n} y_i \, \ell_i(x)
\]
est \textbf{l'unique} polynôme de degré $\leq n$ passant par les $n+1$ points.

\end{mdframed}

\newpage

% ====================================================================
\section*{Exercice 1 --- Construction guidée des polynômes de Lagrange}
% ====================================================================

On souhaite approcher la fonction $\sin$ par un polynôme
passant par les points d'abscisse $0$, $1$ et $2$.

On pose $x_0=0$, $x_1=1$, $x_2=2$.

\begin{enumerate}[itemsep=0.5em]

% ---------- Construction de l_0 ----------
\item On cherche un polynôme $\ell_0$ de degré $2$ vérifiant
\[
  \ell_0(0) = 1, \qquad \ell_0(1) = 0, \qquad \ell_0(2) = 0.
\]
\begin{enumerate}
  \item Puisque $\ell_0$ s'annule en $1$ et en $2$,
        justifier qu'il existe une constante $a$ telle que
        $$\ell_0(x) = a\,(x-1)(x-2)$$

  \item Déterminer $a$ à l'aide de la condition $\ell_0(0) = 1$.
\end{enumerate}

\ifthenelse{\boolean{showSolutions}}{
  \begin{solution}
    $\ell_0$ est de degré $2$ et admet $1$ et $2$ pour racines,
    donc $$\ell_0(x) = a\,(x-1)(x-2)$$
    En évaluant en $0$ :
    \[
      a\,(0-1)(0-2) = 2a = 1, \qquad \text{d'où } a = \frac{1}{2}.
    \]
    Ainsi
    \[
      \ell_0(x) = \frac{(x-1)(x-2)}{2}.
    \]
  \end{solution}
}{}

% ---------- Construction de l_1 et l_2 ----------
\item Par le même raisonnement, construire $\ell_1$ et $\ell_2$ tels que
\[
  \ell_1(0) = 0,\ \ell_1(1) = 1,\ \ell_1(2) = 0
  \qquad \text{et} \qquad
  \ell_2(0) = 0,\ \ell_2(1) = 0,\ \ell_2(2) = 1.
\]

\ifthenelse{\boolean{showSolutions}}{
  \begin{solution}
    $\ell_1$ s'annule en $0$ et $2$, donc $\ell_1(x) = b\,x(x-2)$.
    La condition $\ell_1(1) = 1$ donne $b \cdot 1 \cdot (-1) = 1$,
    soit $b = -1$ :
    \[
      \ell_1(x) = -x(x-2).
    \]
    $\ell_2$ s'annule en $0$ et $1$, donc $\ell_2(x) = c\,x(x-1)$.
    La condition $\ell_2(2) = 1$ donne $c \cdot 2 \cdot 1 = 1$,
    soit $c = \tfrac{1}{2}$ :
    \[
      \ell_2(x) = \frac{x(x-1)}{2}.
    \]
  \end{solution}
}{}

% ---------- Assemblage du polynôme ----------
\item A l'aide de ces polynômes élémentaires $\ell_0$, $\ell_1$ et $\ell_2$, écrire le polynôme $P$ qui passe par les points $(0, \sin 0)$, $(1, \sin 1)$, $(2, \sin 2)$.

\ifthenelse{\boolean{showSolutions}}{
  \begin{solution}
    \[
      P(x) = \sin(0)\,\ell_0(x) + \sin(1)\,\ell_1(x)
            + \sin(2)\,\ell_2(x)
            = \sin(1)\bigl[-x(x-2)\bigr]
            + \sin(2)\,\frac{x(x-1)}{2}.
    \]
  \end{solution}
}{}

% ---------- Lien avec la formule générale ----------
\item Vérifier que $\ell_0$, $\ell_1$, $\ell_2$ coïncident
avec la formule générale donnée dans les rappels :
\[
  \ell_i(x)
  = \prod_{\substack{0 \leq j \leq 2 \\ j \neq i}}
    \frac{x - x_j}{x_i - x_j}.
\]

\ifthenelse{\boolean{showSolutions}}{
  \begin{solution}
    Pour $i = 0$ :
    $\dfrac{(x-x_1)(x-x_2)}{(x_0-x_1)(x_0-x_2)}
    = \dfrac{(x-1)(x-2)}{(-1)(-2)}
    = \dfrac{(x-1)(x-2)}{2} = \ell_0(x)$.
    On retrouve de même les expressions de $\ell_1$ et $\ell_2$.
  \end{solution}
}{}

% ---------- Graphique ----------
\item Tracer à la main sur un même graphique le polynôme $P$ obtenu précédemment
et la courbe de $\sin$ sur $[0,3]$.
Que constate-t-on en dehors de $[0,2]$ ?

\ifthenelse{\boolean{showSolutions}}{
  \begin{solution}
    \begin{center}
    \begin{tikzpicture}[scale=1.1]
      % Shading for interpolation zone
      \fill[ThemeLight!40] (0,-2.6) rectangle (2,1.15);
      \node[Theme!50!black, font=\scriptsize\itshape] at (1,1.05) {interpolation};

      % Shading for extrapolation zone
      \fill[red!5] (2,-2.6) rectangle (4.8,1.15);
      \node[red!40!black, font=\scriptsize\itshape] at (3.5,1.05) {extrapolation};

      % Axes
      \draw[->] (-0.3,0) -- (5.0,0) node[right] {\small$x$};
      \draw[->] (0,-2.7) -- (0,1.35) node[above] {\small$y$};

      % x-tick labels
      \foreach \x in {1,2,3,4}
        \draw (\x,-0.06) -- (\x,0.06) node[below,yshift=-1,font=\scriptsize] {$\x$};

      % y-tick labels
      \foreach \y in {-2,-1,1}
        \draw (-0.06,\y) -- (0.06,\y) node[left,xshift=-1,font=\scriptsize] {$\y$};

      % Plot sin(x)
      \draw[thick,Theme,domain=0:4.7,samples=120,smooth]
        plot (\x, {sin(\x r)});
      \node[Theme!90!black,font=\small,anchor=west] at (4.75,{sin(4.7 r)}) {$\sin x$};

      % Plot P(x)
      \draw[thick,red!80!black,dashed,domain=0:4.7,samples=120,smooth]
        plot (\x, {sin(1 r)*(-\x*(\x-2)) + sin(2 r)*(\x*(\x-1)/2)});
      \node[red!80!black,font=\small,anchor=east] at (4.3,{sin(1 r)*(-4.3*(4.3-2)) + sin(2 r)*(4.3*(4.3-1)/2)}) {$P(x)$};

      % Boundary line between zones
      \draw[gray,densely dotted] (2,-2.6) -- (2,1.15);

      % Mark interpolation points
      \foreach \xi in {0,1,2}
        \filldraw[black] (\xi, {sin(\xi r)}) circle (1.3pt);

      % Dotted vertical drop-lines
      \foreach \xi in {1,2}
        \draw[dotted,gray] (\xi,0) -- (\xi,{sin(\xi r)});
    \end{tikzpicture}
    \end{center}
    Sur $[0,2]$ les deux courbes sont proches.
    Au-delà de $2$, $P$ s'éloigne rapidement de $\sin$ :
    c'est de l'extrapolation, le polynôme n'est plus
    contraint par des données.
  \end{solution}
}{}

% ---------- Comparaison Taylor ----------
\item Comparer avec le polynôme de Taylor de $\sin$
en $0$ au degré $3$ : $T_3(x) = x - \dfrac{x^3}{6}$.
\begin{enumerate}
  \item Lequel approche le mieux $\sin(1{,}5)$ ?
  \item Et $\sin(3)$ ?
\end{enumerate}

\ifthenelse{\boolean{showSolutions}}{
  \begin{solution}
    $T_3(1{,}5) = 1{,}5 - \tfrac{3{,}375}{6} \approx 0{,}9375$.
    \quad $P(1{,}5) \approx 0{,}997$.
    \quad $\sin(1{,}5) \approx 0{,}9975$.

    Lagrange est nettement meilleur car $1{,}5$ est à l'intérieur
    de l'intervalle d'interpolation $[0,2]$.

    \medskip

    Pour $\sin(3) \approx 0{,}141$ :
    $T_3(3) = 3 - 4{,}5 = -1{,}5$ et
    $P(3) = -\sin(1) \cdot 3 + 3\sin(2) \approx 0{,}2$.
    Les deux sont mauvais si loin de leurs zones de validité.
  \end{solution}
}{}

\end{enumerate}

% ====================================================================
\section*{Exercice 2 --- Profil de façade et interpolation}
% ====================================================================

On relève la hauteur $h(x)$, en mètres,
du bord supérieur d'une façade sur une largeur de $6$ mètres.
Les mesures suivantes sont disponibles :
\[
\begin{array}{c|ccc}
  x      & 0 & 3 & 6 \\
  \hline
  h(x)   & 4 & 5 & 3
\end{array}
\]

\begin{enumerate}
\item Déterminer le polynôme de Lagrange $P_2$
      de degré au plus $2$ qui interpole ces trois points.

\ifthenelse{\boolean{showSolutions}}{
  \begin{solution}
    Les polynômes de base sont
    \[
      \ell_0(x)=\frac{(x-3)(x-6)}{(0-3)(0-6)}=\frac{(x-3)(x-6)}{18},
    \]
    \[
      \ell_1(x)=\frac{x(x-6)}{3(3-6)}=-\frac{x(x-6)}{9},
      \qquad
      \ell_2(x)=\frac{x(x-3)}{6 \cdot 3}=\frac{x(x-3)}{18}.
    \]
    Donc
    \[
      P_2(x)=4\ell_0(x)+5\ell_1(x)+3\ell_2(x)
      =-\frac{1}{6}x^2+\frac{5}{6}x+4.
    \]
  \end{solution}
}{}

\item Donner une estimation de la hauteur au point $x=2$.

\ifthenelse{\boolean{showSolutions}}{
  \begin{solution}
    \[
      P_2(2)=-\frac{4}{6}+\frac{10}{6}+4=5\text{ m}.
    \]
  \end{solution}
}{}

\item Donner une estimation de la hauteur au point $x=8$.
      Ce calcul vous semble-t-il fiable ?

\ifthenelse{\boolean{showSolutions}}{
  \begin{solution}
    \[
      P_2(8)=-\frac{64}{6}+\frac{40}{6}+4=0.
    \]
    Le résultat est très fragile : $x=8$ est en dehors
    de l'intervalle de mesures $[0,6]$.
    C'est de l'extrapolation.
  \end{solution}
}{}
\end{enumerate}

% ====================================================================
%  RAPPELS DE COURS : INTÉGRATION NUMÉRIQUE
% ====================================================================
\newpage

\begin{mdframed}[backgroundcolor=ThemeLight, linecolor=Theme,
  linewidth=1.5pt, innertopmargin=10pt, innerbottommargin=10pt]
\textbf{\large Intégration numérique}

\medskip
On rappelle que les ordinateurs ne peuvent calculer que des sommes et des produits. Les polynômes de lagrange peuvent servir à obtenir des valeurs approchées pour des intégrales.

On veut approcher l'intégrale
$\displaystyle I = \int_a^b f(x)\,dx$
à l'aide de $n+1$ points régulièrement espacés
$x_0 = a < x_1 < \cdots < x_n = b$, de pas
$h = \dfrac{b-a}{n}$.

\medskip

\textbf{Rectangles à gauche.}
\[
  R_g = h \sum_{k=0}^{n-1} f(x_k).
\]

\textbf{Trapèzes.}
\[
  T = \frac{h}{2}\bigl(f(x_0) + 2f(x_1)
    + \cdots + 2f(x_{n-1}) + f(x_n)\bigr).
\]

\textbf{Simpson} ($n$ pair).
\[
  S = \frac{h}{3}\bigl(f(x_0) + 4f(x_1) + 2f(x_2)
    + 4f(x_3) + \cdots + 4f(x_{n-1}) + f(x_n)\bigr).
\]

\medskip

\textbf{Idée.}
Chaque méthode approche $f$ par un polynôme de degré
croissant sur chaque sous-intervalle :
constante (rectangles), affine (trapèzes), parabolique (Simpson).
Plus le degré est élevé, meilleure est la précision.
\end{mdframed}

% ====================================================================
\section*{Exercice 3 --- Aire d'une surface à peindre}
% ====================================================================

On dispose des mesures suivantes du profil supérieur
d'une bande de façade :
\[
\begin{array}{c|ccccc}
  x \text{ (m)}    & 0 & 1 & 2 & 3 & 4 \\
  \hline
  f(x) \text{ (m)} & 2{,}0 & 2{,}4 & 2{,}8 & 2{,}5 & 2{,}2
\end{array}
\]
On cherche à approcher l'aire de la bande comprise
entre la courbe $y=f(x)$ et l'axe horizontal.

Pour chaque méthode ci-dessous, représenter sur un dessin la zone dont on cherche à calculer l'aire et l'approximation obtenue.

\begin{enumerate}
\item Donner l'approximation par la méthode des rectangles
      à gauche avec un pas $h=1$.

\ifthenelse{\boolean{showSolutions}}{
  \begin{solution}
    \[
      R_g = 1 \times \bigl(f(0)+f(1)+f(2)+f(3)\bigr)
      = 2{,}0+2{,}4+2{,}8+2{,}5 = 9{,}7\ \text{m}^2.
    \]
  \end{solution}
}{}

\item Donner l'approximation par la méthode des trapèzes.

\ifthenelse{\boolean{showSolutions}}{
  \begin{solution}
    \[
      T = \frac{1}{2}\bigl(f(0)+2f(1)+2f(2)+2f(3)+f(4)\bigr)
      = \frac{1}{2}(2{,}0+4{,}8+5{,}6+5{,}0+2{,}2) = 9{,}8\ \text{m}^2.
    \]
  \end{solution}
}{}

\item Donner l'approximation par la méthode de Simpson.

\ifthenelse{\boolean{showSolutions}}{
  \begin{solution}
    Le nombre de sous-intervalles vaut $4$ (pair),
    on peut utiliser Simpson :
    \[
      S = \frac{1}{3}\bigl(f(0)+4f(1)+2f(2)+4f(3)+f(4)\bigr)
      = \frac{1}{3}(2{,}0+9{,}6+5{,}6+10{,}0+2{,}2)
      \approx 9{,}8\ \text{m}^2.
    \]
  \end{solution}
}{}

\item Si un pot de peinture couvre $12\ \text{m}^2$,
      combien de pots faut-il prévoir pour deux couches ?

\ifthenelse{\boolean{showSolutions}}{
  \begin{solution}
    Avec Simpson l'aire est d'environ $9{,}8\ \text{m}^2$
    pour une couche.
    Pour deux couches :
    $2 \times 9{,}8 = 19{,}6\ \text{m}^2$.
    Il faut donc prévoir $2$ pots.
  \end{solution}
}{}
\end{enumerate}

\newpage

% ====================================================================
\section*{Exercice 4 --- Mesures sur éprouvettes}
% ====================================================================

Lors d'un essai, on mesure l'allongement $u$ en fonction
de la force appliquée $F$ :
\[
\begin{array}{c|ccc}
  F \text{ (kN)}  & 0 & 10 & 20 \\
  \hline
  u \text{ (mm)}  & 0 & 1{,}2 & 2{,}1
\end{array}
\]

\begin{enumerate}
\item Déterminer le polynôme d'interpolation de degré
      au plus $2$ reliant $u$ à $F$.

\item En déduire une estimation de l'allongement pour $F=15$ kN.

\item Comparer cette approche avec une simple interpolation
      affine entre $F=10$ et $F=20$.

\ifthenelse{\boolean{showSolutions}}{
  \begin{solution}
    On cherche $P(F) = aF^2 + bF + c$.
    Comme $P(0) = 0$, on a $c = 0$.
    Puis
    \[
      100a + 10b = 1{,}2, \qquad 400a + 20b = 2{,}1.
    \]
    En divisant la seconde par $2$ et en soustrayant :
    $100a = -0{,}15$, soit $a = -0{,}0015$.
    Puis $b = 0{,}135$.
    \[
      P(F) = -0{,}0015\,F^2 + 0{,}135\,F.
    \]
    Pour $F = 15$ :
    \[
      P(15) = -0{,}0015 \times 225 + 0{,}135 \times 15
             = 1{,}6875\ \text{mm}.
    \]
    L'interpolation affine entre $10$ et $20$ donne
    \[
      1{,}2 + \frac{15-10}{20-10}(2{,}1 - 1{,}2) = 1{,}65\ \text{mm}.
    \]
    Les deux valeurs sont proches.
    L'interpolation quadratique tient compte de la courbure
    du comportement, tandis que l'interpolation affine suppose
    une relation linéaire entre les deux points.
  \end{solution}
}{}
\end{enumerate}

\vspace{2em}
% ====================================================================
\section*{Exercice 5 --- Erreur de quadrature}
% ====================================================================

On considère l'intégrale
\[
  I = \int_0^1 e^x\,dx = e - 1.
\]

\begin{enumerate}
\item Donner l'approximation de $I$ par la méthode des trapèzes
      avec un seul sous-intervalle.

\item Donner l'approximation de $I$ par la méthode de Simpson
      avec deux sous-intervalles.

\item Comparer les erreurs absolues obtenues.

\ifthenelse{\boolean{showSolutions}}{
  \begin{solution}
    $I = e - 1 \approx 1{,}7183$.

    \medskip

    Trapèzes :
    $T = \tfrac{1}{2}(f(0)+f(1)) = \tfrac{1+e}{2} \approx 1{,}8591$.
    Erreur $\approx 0{,}14$.

    \medskip

    Simpson :
    $S = \tfrac{1}{6}\bigl(f(0)+4f(\tfrac{1}{2})+f(1)\bigr)
    = \tfrac{1}{6}(1+4e^{1/2}+e) \approx 1{,}7189$.
    Erreur $\approx 0{,}0006$.

    \medskip

    Simpson est nettement plus précis sur cet exemple :
    l'erreur est environ $200$ fois plus petite.
  \end{solution}
}{}
\end{enumerate}

% ====================================================================
\section*{Exercice 6 --- Retrouver la formule de Simpson}
% ====================================================================

On considère une fonction $f$ continue sur un intervalle $[a,b]$.
On pose $m = \dfrac{a+b}{2}$ le milieu de l'intervalle
et $h = \dfrac{b-a}{2}$.

L'idée est d'approcher $f$ par le polynôme de Lagrange de degré $2$
passant par les trois points $(a, f(a))$, $(m, f(m))$, $(b, f(b))$,
puis d'intégrer ce polynôme à la place de $f$.

\begin{enumerate}

\item Écrire les trois polynômes de base de Lagrange
$\ell_0$, $\ell_1$, $\ell_2$ associés aux n\oe uds
$x_0 = a$, $x_1 = m$, $x_2 = b$.

\ifthenelse{\boolean{showSolutions}}{
  \begin{solution}
    \[
      \ell_0(x) = \frac{(x-m)(x-b)}{(a-m)(a-b)}
               = \frac{(x-m)(x-b)}{2h^2}\,,
    \]
    \[
      \ell_1(x) = \frac{(x-a)(x-b)}{(m-a)(m-b)}
               = -\frac{(x-a)(x-b)}{h^2}\,,
    \]
    \[
      \ell_2(x) = \frac{(x-a)(x-m)}{(b-a)(b-m)}
               = \frac{(x-a)(x-m)}{2h^2}\,.
    \]
  \end{solution}
}{}

\item On effectue le changement de variable $t = \dfrac{x-a}{h}$
(de sorte que $t = 0$ en $x = a$, $t = 1$ en $x = m$,
$t = 2$ en $x = b$).
Montrer que
\[
  \ell_0(x) = \frac{(t-1)(t-2)}{2}\,,
  \qquad
  \ell_1(x) = -t(t-2),
  \qquad
  \ell_2(x) = \frac{t(t-1)}{2}\,.
\]

\ifthenelse{\boolean{showSolutions}}{
  \begin{solution}
    Avec $x = a + th$, on a $x - a = th$, $x - m = (t-1)h$,
    $x - b = (t-2)h$.
    En substituant :
    \[
      \ell_0 = \frac{(t-1)h \cdot (t-2)h}{2h^2}
             = \frac{(t-1)(t-2)}{2}\,.
    \]
    De même pour $\ell_1$ et $\ell_2$.
  \end{solution}
}{}

\item Calculer les trois intégrales
$\displaystyle\int_a^b \ell_i(x)\,dx$
en utilisant le changement de variable
($dx = h\,dt$, bornes $0$ et $2$).

\ifthenelse{\boolean{showSolutions}}{
  \begin{solution}
    \[
      \int_a^b \ell_0(x)\,dx
      = h\int_0^2 \frac{(t-1)(t-2)}{2}\,dt
      = \frac{h}{2}\Bigl[\frac{t^3}{3}-\frac{3t^2}{2}+2t\Bigr]_0^2
      = \frac{h}{2}\times\frac{2}{3}
      = \frac{h}{3}\,.
    \]
    \[
      \int_a^b \ell_1(x)\,dx
      = h\int_0^2 \bigl(-t^2+2t\bigr)\,dt
      = h\Bigl[-\frac{t^3}{3}+t^2\Bigr]_0^2
      = h\times\frac{4}{3}
      = \frac{4h}{3}\,.
    \]
    \[
      \int_a^b \ell_2(x)\,dx
      = h\int_0^2 \frac{t(t-1)}{2}\,dt
      = \frac{h}{2}\Bigl[\frac{t^3}{3}-\frac{t^2}{2}\Bigr]_0^2
      = \frac{h}{2}\times\frac{2}{3}
      = \frac{h}{3}\,.
    \]
  \end{solution}
}{}

\item En déduire que
\[
  \int_a^b f(x)\,dx
  \;\approx\;
  \int_a^b P(x)\,dx
  \;=\;
  \frac{b-a}{6}
  \left[f(a) + 4\,f\!\left(\frac{a+b}{2}\right) + f(b)\right].
\]
Reconnaître la formule de Simpson sur un intervalle.

\ifthenelse{\boolean{showSolutions}}{
  \begin{solution}
    \[
      \int_a^b P(x)\,dx
      = f(a)\cdot\frac{h}{3}
        + f(m)\cdot\frac{4h}{3}
        + f(b)\cdot\frac{h}{3}
      = \frac{h}{3}\bigl[f(a) + 4f(m) + f(b)\bigr].
    \]
    Comme $h = \dfrac{b-a}{2}$, on obtient
    $\dfrac{h}{3} = \dfrac{b-a}{6}$, d'où
    \[
      \int_a^b P(x)\,dx
      = \frac{b-a}{6}\bigl[f(a) + 4f(m) + f(b)\bigr].
    \]
    C'est exactement la formule de Simpson
    sur l'intervalle $[a,b]$.
  \end{solution}
}{}

\end{enumerate}

\end{document}
